\section{Entrenamiento de MDLM}

\subsection{Función de pérdida ELBO}

En los modelos AR, el objetivo del entrenamiento es maximizar directamente la verosimilitud logarítmica $\log p(\mathbf{x})$. En los modelos de difusión no es posible calcular con precisión la verosimilitud, por lo que se utiliza el \textbf{límite inferior de evidencia} (Evidence Lower Bound, ELBO) como objetivo del entrenamiento.

Una de las principales contribuciones de MDLM es haber derivado un ELBO en tiempo continuo \textbf{extremadamente conciso}:

$$\mathcal{L}_{\text{MDLM}} = \mathbb{E}_{t \sim \mathcal{U}(0,1), \, \mathbf{z}_t \sim q(\mathbf{z}_t|\mathbf{x})} \left[ \frac{\alpha_t'}{1 - \alpha_t} \sum_{i: z_t^{(i)} = \text{M}} -\log x_\theta^{(i)}(\mathbf{z}_t, t) \right]$$

Significado de cada término:
\begin{itemize}
    \item $\mathbb{E}_{t \sim \mathcal{U}(0,1)}$: Se muestra el paso de tiempo $t$ uniformemente de $[0, 1]$
    \item $\mathbf{z}_t \sim q(\mathbf{z}_t | \mathbf{x})$: Se añade ruido al texto limpio $\mathbf{x}$ para obtener $\mathbf{z}_t$
    \item $\frac{\alpha_t'}{1 - \alpha_t}$: \textbf{Término de ponderación}, es decir, la probabilidad de desmascaramiento (chance of unmasking)
    \item $-\log x_\theta^{(i)}(\mathbf{z}_t, t)$: \textbf{Pérdida de entropía cruzada}, calculada únicamente en las posiciones con máscara
\end{itemize}

\begin{importantbox}{Límite continuo}
La idea clave del artículo sobre MDLM es que, cuando el número de pasos de discretización $T \to \infty$, la cadena de Markov de tiempo discreto converge a una cadena de Markov de tiempo continuo, obteniendo así la función de pérdida concisa anterior. Esta simplificación facilita enormemente el entrenamiento y la evaluación.
\end{importantbox}

\subsection{Comparación entre el entrenamiento AR y MDLM}

Los flujos de entrenamiento de los modelos AR y MDLM son muy similares; solo existen tres diferencias clave:

\begin{table}[H]
\centering
\caption{Comparación del entrenamiento de modelos AR y MDLM}
\begin{tabular}{@{}lll@{}}
\toprule
\textbf{Aspecto} & \textbf{Modelo AR} & \textbf{MDLM} \\
\midrule
Mecanismo de atención & Atención causal & Atención bidireccional \\
Entrada & Texto limpio $\mathbf{x}$ & Secuencia con ruido $\mathbf{z}_t$ (después de máscaras aleatorias) \\
Muestreo adicional & Ninguno & Se requiere muestrear el paso de tiempo $t$ \\
Función de pérdida & Entropía cruzada & Entropía cruzada ponderada (peso $\frac{\alpha_t'}{1-\alpha_t}$) \\
Posición de cálculo & Todas las posiciones & Solo las posiciones con máscara \\
\bottomrule
\end{tabular}
\end{table}

\begin{knowledgebox}{Concisión del entrenamiento}
Si sabes cómo entrenar un modelo AR, básicamente ya sabes cómo entrenar uno MDLM. La única operación adicional es: (1) Muestrear aleatoriamente un paso de tiempo $t$; (2) Aplicar máscaras aleatorias a la entrada según $t$; (3) Incluir el término de ponderación en la pérdida. Esta concisión es una razón importante por la que MDLM ha sido ampliamente adoptado.
\end{knowledgebox}

\subsection{Detalles del proceso de muestreo}

El proceso completo de muestreo de MDLM (muestreo ancestral, Ancestral Sampling):

\begin{enumerate}
    \item \textbf{Inicialización}: $\mathbf{z}_T = [\text{M}, \text{M}, \ldots, \text{M}]$ (secuencia completamente enmascarada)
    \item \textbf{Denoising de un paso} (de $\mathbf{z}_t$ a $\mathbf{z}_s$):
    \begin{enumerate}
        \item Enviar $\mathbf{z}_t$ a un Transformer bidireccional
        \item El modelo predice la distribución de palabras para todas las posiciones con máscara
        \item Muestrear de la distribución para rellenar todas las posiciones con máscara
        \item Decidir qué predicciones conservar según la probabilidad de desmascaramiento $\frac{\alpha_s - \alpha_t}{1 - \alpha_t}$
        \item Las predicciones no conservadas se vuelven a enmascarar
        \item Los tokens revelados permanecen sin cambios
    \end{enumerate}
    \item Repetir el paso 2 un total de $T$ veces hasta que $\mathbf{z}_0$ esté completamente revelado
\end{enumerate}

\begin{warningbox}{Esto no es un método heurístico}
Este proceso de muestreo no es un método heurístico diseñado por intuición, sino un muestreador ancestral (Ancestral Sampler) derivado estrictamente de deducciones matemáticas. Las deducciones matemáticas requieren que: (1) Los tokens revelados deben permanecer sin cambios; (2) Solo las predicciones en posiciones con máscara se vuelven a enmascarar.
\end{warningbox}

\subsection{Resumen de la sección}

La función de pérdida de entrenamiento de MDLM es una entropía cruzada ponderada, donde el peso está determinado por la probabilidad de desmascaramiento. El límite continuo es la simplificación central de MDLM en comparación con trabajos anteriores. El flujo de entrenamiento es muy similar al de los modelos AR, lo que reduce la barrera para su implementación.

\newpage
